\chapter*{Введение}
\addcontentsline{toc}{chapter}{Введение}
\label{ch:intro}

Современные технологии стремительно меняют способы, которыми люди взаимодействуют с информацией, включая контент, такой как фильмы и сериалы. В эпоху, когда объем доступного контента непрерывно растет, пользователям становится всё сложнее ориентироваться в этом разнообразии и находить произведения, которые соответствуют их вкусам и предпочтениям. Традиционные методы поиска уже не способны справляться с такой нагрузкой, и на первый план выходят рекомендательные системы, которые помогают автоматизировать процесс выбора контента. Такие системы позволяют не только экономить время пользователей, но и открывать им новые возможности для открытия уникальных и интересных фильмов, которые они могли бы упустить. Они применяются в потоковых сервисах, онлайн-кинотеатрах и многих других платформах, становясь неотъемлемой частью цифрового мира.

Проект, представленный в данном документе, направлен на создание платформы, которая объединяет возможности хранения, обработки данных и персонализации рекомендаций. Он предлагает пользователям удобный способ поиска фильмов, а также предоставляет функционал для взаимодействия с контентом, включая возможность просматривать информацию о фильмах, оценивать их и получать персонализированные рекомендации. Система включает в себя такие функции, как регистрация и аутентификация пользователей, защищённый доступ к данным, управление контентом через админ-панель, а также использование продвинутых алгоритмов машинного обучения для формирования рекомендаций. Таким образом, проект создаёт основу для удобного и безопасного взаимодействия с фильмами, ориентированного на предпочтения каждого пользователя.

\section*{Описание проекта}

Наша команда разработала платформу для хранения и обработки данных о фильмах, а также предоставления пользователям персонализированных рекомендаций. Проект состоит из трёх ключевых компонентов:

\begin{itemize}
    \item Бэкенд для управления данными о фильмах, тегах и связанных сущностях.
    \item Модуль машинного обучения, который анализирует данные и генерирует рекомендации.
    \item Фронтенд, обеспечивающий удобный интерфейс для пользователей, включая поисковую строку, отображение рекомендованных фильмов и просмотр информации о фильмах.
\end{itemize}

В проекте реализованы следующие функциональные возможности:
\begin{itemize}
    \item Персонализированные рекомендации на основе алгоритмов машинного обучения, таких как UserKNN и TF-IDF, что позволяет учитывать как вкусы пользователя, так и сходство контента.
    \item Регистрация и аутентификация пользователей, что обеспечивает индивидуальный опыт работы с системой и защиту личных данных.
    \item Админ-панель для добавления и управления контентом, включая фильмы, жанры, теги и режиссёров, что упрощает администрирование и обновление базы данных.
    \item Возможность просмотра данных о фильмах, включая описание, жанры, теги, режиссёров и комментарии, что делает платформу полноценным источником информации о кино.
    \item Функция оценки фильмов пользователями с помощью системы рейтингов, что позволяет собирать обратную связь и улучшать точность рекомендаций.
    \item Рекомендации фильмов на основе поставленных пользователем оценок, что даёт более персонализированный подход к выбору контента.
    \item Рекомендации фильмов на основе текущего фильма, который просматривает пользователь, что добавляет контекстную релевантность предложений.
\end{itemize}

Наряду с этим, проект удовлетворяет важным нефункциональным требованиям:
\begin{itemize}
    \item Доступ к личным данным пользователей защищён авторизацией и аутентификацией, что минимизирует риски утечки информации.
    \item Все передачи данных осуществляются через протокол HTTPS, обеспечивая безопасность и конфиденциальность.
    \item Архитектура проекта спроектирована так, чтобы быть масштабируемой и поддерживать рост числа пользователей и объёма данных.
    \item Развёртывание системы реализовано таким образом, чтобы быть простым и удобным для разработчиков.
\end{itemize}

В этом документе подробно рассматриваются структура системы, её основные компоненты, функциональные и нефункциональные требования.

\section*{Бэкенд}

Бэкенд нашего приложения разработан на базе фреймворка Django с использованием Django REST Framework (DRF). Он выполняет следующие ключевые функции:
\begin{itemize}
    \item Хранение информации о фильмах, тегах и других метаданных.
    \item Предоставление API для выполнения CRUD операций:
    \begin{itemize}
        \item Создание новых записей о фильмах и связанных данных.
        \item Чтение существующих данных.
        \item Обновление информации, например, изменение тегов или других атрибутов фильмов.
        \item Удаление записей из базы данных.
    \end{itemize}
    \item Управление пользователями через регистрацию, аутентификацию и авторизацию, что обеспечивает персонализацию работы с системой.
\end{itemize}

Для удобства администраторов предоставлена удобная админ-панель, через которую можно управлять контентом платформы. Она позволяет добавлять и обновлять информацию о фильмах, включая жанры, теги и режиссёров, а также модерировать пользовательский контент, например, комментарии.

Архитектура бэкенда организована с учётом масштабируемости и интеграции с другими компонентами системы, такими как модуль машинного обучения и фронтенд. Использование Django и DRF обеспечивает:
\begin{itemize}
    \item Высокую скорость разработки благодаря удобным инструментам.
    \item Структурированность и поддержку принципов MVC.
    \item Гибкость в настройке API и интеграции с внешними системами.
    \item Поддержку современных стандартов безопасности, включая использование HTTPS для всех передаваемых данных и защиты личных данных пользователей.
\end{itemize}

\section*{Модуль машинного обучения}

Модуль машинного обучения является ядром нашей рекомендательной системы. Он реализует два подхода к созданию рекомендаций:
\begin{itemize}
    \item User-based подход (UserKNN):
    \begin{itemize}
        \item Использует алгоритм ближайших соседей (k-nearest neighbors) для поиска пользователей с похожими предпочтениями.
        \item Основывается на анализе пользовательских взаимодействий с фильмами (оценки, просмотры и т.д.).
        \item Предоставляет рекомендации, исходя из данных о схожих пользователях, что делает систему более персонализированной.
    \end{itemize}
    \item Content-based подход:
    \begin{itemize}
        \item Использует метод TF-IDF для анализа текстового содержимого (теги, описания фильмов и т.д.).
        \item Вычисляет косинусное сходство между фильмами для определения степени их схожести.
        \item Рекомендует фильмы, которые максимально соответствуют интересам пользователя на основе их взаимодействия с похожим контентом.
    \end{itemize}
\end{itemize}

Модуль обновляет данные о рекомендациях раз в сутки, загружая их в базу данных через специальный API. Такой подход позволяет сохранять высокую производительность системы и гарантировать, что рекомендации всегда актуальны. Благодаря этому пользователи получают рекомендации как на основе своих оценок, так и на основе текущего фильма, который они просматривают. Это делает рекомендации не только точными, но и контекстно релевантными.

\section*{Фронтенд}

Фронтенд является пользовательским интерфейсом системы, предоставляющим интуитивно понятный способ поиска, просмотра и получения рекомендаций фильмов. Основные функции фронтенда:
\begin{itemize}
    \item Поисковая строка:
    \begin{itemize}
        \item Позволяет пользователям быстро находить фильмы по названию или ключевым словам.
        \item Отображает персонализированные рекомендации на основе данных из модуля машинного обучения.
    \end{itemize}
    \item Отображение фильмов:
    \begin{itemize}
        \item Информация о фильме, включая описание, постер, теги и рекомендации.
        \item Подробные данные о режиссёрах, жанрах и комментариях.
    \end{itemize}
    \item Просмотр фильмов:
    \begin{itemize}
        \item Пользователи могут просматривать карточки фильмов с подробной информацией, включая оценки, отзывы и связанные теги.
        \item Доступ к списку рекомендованных фильмов, основанных на предпочтениях пользователя или сходстве с текущим фильмом.
    \end{itemize}
\end{itemize}

Фронтенд взаимодействует с бэкендом через REST API, что обеспечивает гибкость и масштабируемость системы. Все передачи данных защищены использованием протокола HTTPS, что обеспечивает безопасность пользовательской информации.

\section*{Цели проекта}

Основная цель проекта — создание эффективной системы рекомендаций, которая:
\begin{itemize}
    \item Помогает пользователям находить фильмы, соответствующие их вкусам и предпочтениям.
    \item Упрощает взаимодействие с большим объемом данных о фильмах.
    \item Использует современные технологии для персонализации рекомендаций и повышения удовлетворенности пользователей.
    \item Расширяет горизонты пользователей, предлагая им фильмы, которые они могли бы упустить из-за ограниченности традиционных методов поиска.
    \item Обеспечивает безопасность пользовательских данных через защиту доступа и передачу данных.
    \item Гарантирует удобство администрирования, предоставляя мощные инструменты для управления контентом.
\end{itemize}

Таким образом, наша система способствует развитию культуры потребления качественного контента и упрощает доступ к разнообразным произведениям киноискусства. Мы стремимся предоставить пользователям платформу, которая объединяет удобство, персонализацию и инновации, делая их взаимодействие с системой максимально комфортным и эффективным.

Этот проект не только упрощает процесс поиска фильмов, но и служит инструментом для изучения пользовательских предпочтений, что открывает возможности для дальнейших исследований в области анализа данных и улучшения рекомендательных алгоритмов. Мы уверены, что внедрение подобных технологий не только способствует повышению пользовательского опыта, но и задаёт новый стандарт для современных цифровых платформ.

Данный проект представляет собой сочетание современных технологий разработки (Django, DRF) и машинного обучения (UserKNN, TF-IDF), что позволяет реализовать высокоэффективную платформу для рекомендаций фильмов. С помощью этой системы пользователи могут значительно упростить процесс поиска фильмов, обнаружить новые шедевры и получить персонализированный опыт. Наши алгоритмы создают мост между предпочтениями пользователей и огромным массивом доступного контента, делая выбор фильмов простым и увлекательным.

Кроме того, мы уверены, что данная платформа станет основой для дальнейших улучшений и дополнений, включая интеграцию с новыми источниками данных, разработку более сложных моделей и расширение функционала для пользователей. В следующих разделах документации будут детально рассмотрены архитектура системы, принципы работы модулей и примеры использования, что поможет лучше понять её потенциал и возможности.
\endinput
